\section{Locally Finite Limit}
The recurrence determines every finite stage. To obtain one continuous
function on an unbounded domain, we must also justify its infinite sum.
Here activation gives something stronger than pointwise convergence: on
each bounded interval, sufficiently late stages agree exactly.

\begin{itemize}
\item Exact stabilization on bounded intervals.
\item Continuity of the completed curve.
\end{itemize}

For real $t\ge2$, define
\begin{equation*}
\begin{aligned}
h(t)&=\sin^2(\pi t)+\sum_{m=2}^{\infty}a_m(t)W_m(t),\\
F(t)&=t+i h(t).
\end{aligned}
\end{equation*}

Fix $X\ge2$. Every term with $m\ge X$ vanishes throughout $[2,X]$, since
$t\le X\le m$ implies $a_m(t)=0$. Thus the infinite expression is a finite
sum on that entire interval. In particular,
\begin{equation*}
\begin{aligned}
M\ge\lceil X\rceil
&\ \Longrightarrow\ h_M(t)=h(t)\quad(2\le t\le X)
&&[\text{Activation}]\\
\sup_{2\le t\le X}|h_M(t)-h(t)|&=0.
&&[\text{By Definition}]
\end{aligned}
\end{equation*}

Each finite stage is continuous. Exact agreement on compact intervals
therefore proves both local uniform convergence and continuity of the limit:
\begin{equation*}
\begin{aligned}
h_M&\longrightarrow h\quad\text{uniformly on every }[2,X],\\
F&\text{ is continuous on }[2,\infty).
\quad\blacksquare\ \text{[Q.E.D.]}
\end{aligned}
\end{equation*}

Local finiteness refers to the number of active filter terms on a bounded
input interval. It does not bound the number of primes represented by the
whole curve. At $t=12$, only indices $2$ through $11$ can contribute; at
$t=12.5$, index $12$ can contribute as well. Increasing the input permits
further filters without changing the function already defined on earlier
bounded intervals.

No global uniform-convergence claim on $[2,\infty)$ is needed. Nor does a
continuous piecewise-linear activation establish differentiability at its
breakpoints. The construction is a continuous complex-valued function of a
real variable; holomorphicity is not asserted.
